\begin{longtable}{|p{1.5cm}|p{5.0cm}|p{5.0cm}|p{3.5cm}|}
\hline
\rowcolor{headercolor}
\textcolor{white}{\textbf{Semaine}} & \textcolor{white}{\textbf{Binôme 1 (Moi)}} & \textcolor{white}{\textbf{Binôme 2 (Collègue)}} & \textcolor{white}{\textbf{Livrables / Jalons}} \\
\hline
\endfirsthead
\hline
\rowcolor{headercolor}
\textcolor{white}{\textbf{Semaine}} & \textcolor{white}{\textbf{Binôme 1 (Moi)}} & \textcolor{white}{\textbf{Binôme 2 (Collègue)}} & \textcolor{white}{\textbf{Livrables / Jalons}} \\
\hline
\endhead
\hline
\endfoot
\hline
\endlastfoot

\textbf{S1} \newline \scriptsize 02/02 -- 08/02 & Étude du contexte de l'orientation universitaire tunisienne, analyse des objectifs de CapAvenir, étude de l'environnement technique, rédaction des premières notes de cadrage. & Collecte des documents de référence, étude de l'état de l'art des solutions d'aide à l'orientation, identification des parties prenantes, synthèse des besoins fonctionnels. & Note de cadrage initial, Rapport de démarrage. \\
\hline
\rowcolor{rowgray}
\textbf{S2} \newline \scriptsize 09/02 -- 15/02 & Analyse fonctionnelle et comparative des plateformes existantes, identification des limites, spécification des opportunités d'amélioration. & Veille technologique sur les algorithmes de recommandation, étude des approches psychométriques (RIASEC) et adaptatives (CAT-IRT). & Rapport d'étude de l'existant. \\
\hline
\textbf{S3} \newline \scriptsize 16/02 -- 22/02 & Recueil et spécification des besoins de l'Espace Étudiant, élaboration des scénarios d'utilisation principaux, rédaction préliminaire du rapport. & Analyse et spécification des besoins des Espaces Conseiller et Administrateur, définition des exigences non fonctionnelles et de sécurité. & Cahier des charges fonctionnel, Premier jet du rapport (Ch. 1). \\
\hline
\rowcolor{rowgray}
\textbf{S4} \newline \scriptsize 23/02 -- 01/03 & Rédaction des cas d'utilisation (tests adaptatifs, recommandations, simulations), élaboration des diagrammes de cas d'utilisation UML. & Modélisation des cas d'utilisation pour les Espaces Conseiller et Administrateur, modélisation des flux métiers généraux. & Spécifications fonctionnelles validées, Diagrammes UML. \\
\hline
\textbf{S5} \newline \scriptsize 02/03 -- 08/03 & Conception de l'architecture globale (MVC + Couche Service), définition de l'organisation modulaire, documentation des choix techniques. & Élaboration des diagrammes de séquence et d'activités UML, modélisation des interactions entre les composants du système. & Dossier de conception logicielle, Diagrammes de séquence UML. \\
\hline
\rowcolor{rowgray}
\textbf{S6} \newline \scriptsize 09/03 -- 15/03 & Conception du modèle conceptuel de données (MCD) et des relations, définition des contraintes d'intégrité référentielle. & Structuration physique du stockage (tests, résultats, historiques), rédaction des scripts de création SQL et des migrations. & Modèle Conceptuel et Physique de Données (MCD/MPD). \\
\hline
\textbf{S7} \newline \scriptsize 16/03 -- 22/03 & Initialisation de l'environnement de développement backend (Laravel), implémentation de l'authentification et de la gestion des rôles (RBAC). & Configuration de la base de données et des migrations, développement des services de base de gestion des utilisateurs. & Socle applicatif backend, Schéma de base de données opérationnel. \\
\hline
\rowcolor{rowgray}
\textbf{S8} \newline \scriptsize 23/03 -- 29/03 & Développement du moteur de test adaptatif (CAT), implémentation de la collecte des réponses, tests unitaires sur les calculs RIASEC. & Développement du moteur de recommandation (SIAEPI v8.0), intégration des règles d'admissibilité basées sur le Score FG. & Moteur d'orientation et de recommandation fonctionnel. \\
\hline
\textbf{S9} \newline \scriptsize 30/03 -- 05/04 & Intégration frontend de l'Espace Étudiant (connexion, profil étudiant, interfaces dynamiques du test adaptatif). & Intégration frontend de l'affichage des recommandations de filières, de la recherche de formations et des favoris. & Interfaces de l'Espace Étudiant opérationnelles. \\
\hline
\rowcolor{rowgray}
\textbf{S10} \newline \scriptsize 06/04 -- 12/04 & Développement des modules de simulation de notes et de simulation de spécialités, intégration des calculs associés. & Développement des outils de comparaison de filières et du simulateur de carrières professionnelles. & Outils d'aide à la décision et simulateurs opérationnels. \\
\hline
\textbf{S11} \newline \scriptsize 13/04 -- 19/04 & Développement du tableau de bord conseiller, implémentation du suivi de cohorte et des indicateurs de progression. & Développement des services de validation des recommandations, du cycle de coaching collaboratif et des notifications. & Espace Conseiller (CRM) opérationnel. \\
\hline
\rowcolor{rowgray}
\textbf{S12} \newline \scriptsize 20/04 -- 26/04 & Développement du module d'administration (gestion des utilisateurs et rôles, sécurité, logs système). & Développement du module de gestion du catalogue de formations et d'importation de données massives (Excel). & Espace Administration opérationnel. \\
\hline
\textbf{S13} \newline \scriptsize 27/04 -- 03/05 & Intégration globale des modules Frontend et Backend, résolution des conflits d'intégration, premiers tests système. & Optimisation des requêtes SQL et indexation, vérification de la cohérence globale des flux, déploiement sur serveur de préproduction. & Version intégrée globale, Plateforme déployée en préproduction. \\
\hline
\rowcolor{rowgray}
\textbf{S14} \newline \scriptsize 04/05 -- 10/05 & Élaboration des plans de tests unitaires et d'intégration, exécution des scénarios de test sur l'Espace Étudiant, correction des bugs. & Exécution des tests fonctionnels et de sécurité sur les Espaces Conseiller et Administrateur, validation avec les encadrants. & Cahier de tests validé, Version stable de l'application. \\
\hline
\textbf{S15} \newline \scriptsize 11/05 -- 17/05 & Optimisation ergonomique de l'interface utilisateur, réduction de la latence, tests de non-régression finaux. & Audit de sécurité de l'application, durcissement du serveur, déploiement final en production. & Application déployée en production. \\
\hline
\rowcolor{rowgray}
\textbf{S16} \newline \scriptsize 18/05 -- 24/05 & Finalisation de la rédaction des chapitres d'introduction, d'analyse des besoins et de conception du rapport. & Rédaction des chapitres de réalisation technique, de tests d'évaluation et de conclusion, mise en page des illustrations. & Rapport de PFE finalisé et prêt pour impression. \\
\hline
\textbf{S17} \newline \scriptsize 25/05 -- 30/05 & Conception du support de présentation (diapositives), préparation de la soutenance orale, simulations de passage. & Préparation et scénarisation de la démonstration technique de l'application, répétitions finales en binôme. & Support de soutenance finalisé, Démo prête. \\
\hline
\end{longtable}
